% Gabarit pandoc des corrigés publiés (brut et guidé). Contenu sans marque ni nom de site.
% Compilation : lualatex (latexmk -lualatex).
% Le Markdown reste la source unique : ne pas éditer ce fichier .tex à la main.
\documentclass[11pt,a4paper]{article}

% --- Moteur, polices (Latin Modern) et langue ---------------------------------
\usepackage{iftex}
\ifLuaTeX\else\errmessage{Ce document se compile avec lualatex}\fi
\usepackage{amsmath,amssymb,mathtools}
\usepackage{fontspec}
% mathrm=sym : \mathrm prend ses glyphes dans la police mathématique (ℓ, μ, Ω), comme KaTeX.
\usepackage[math-style=TeX,bold-style=TeX,mathrm=sym]{unicode-math}
% Repli pour les glyphes absents de Latin Modern (indices et exposants Unicode,
% flèches, filets des schémas en texte, symboles) : DejaVu, puis Latin Modern Math
% (ℓ, lettres mathématiques Unicode des légendes).
\directlua{
  luaotfload.add_fallback("repliTexte", {"DejaVuSerif:mode=harf;", "DejaVuSans:mode=harf;", "file:latinmodern-math.otf:mode=harf;"})
  luaotfload.add_fallback("repliMono", {"DejaVuSansMono:mode=harf;", "DejaVuSans:mode=harf;", "file:latinmodern-math.otf:mode=harf;"})
}
\setmainfont{Latin Modern Roman}[RawFeature={fallback=repliTexte}]
\setsansfont{Latin Modern Sans}[RawFeature={fallback=repliTexte}]
\setmonofont{Latin Modern Mono}[RawFeature={fallback=repliMono}]
\setmathfont{Latin Modern Math}
\usepackage[french]{babel}
\frenchsetup{SmallCapsFigTabCaptions=false}

% --- Mise en page --------------------------------------------------------------
\usepackage[a4paper,margin=2.2cm,headheight=14pt,footskip=1.2cm]{geometry}
\usepackage{microtype}
\setlength{\emergencystretch}{4em}
\tolerance=1500
\setlength{\parindent}{0pt}
\setlength{\parskip}{0.55\baselineskip plus 2pt minus 1pt}
\usepackage{xcolor}
% Gris nommés : babel-french rend « ! » actif dans le document.
\definecolor{gris2}{gray}{0.98}
\definecolor{gris3}{gray}{0.97}
\definecolor{gris5}{gray}{0.95}
\definecolor{gris35}{gray}{0.65}
\definecolor{gris40}{gray}{0.60}
\definecolor{gris45}{gray}{0.55}
\definecolor{gris50}{gray}{0.50}
\definecolor{gris55}{gray}{0.45}
\definecolor{gris60}{gray}{0.40}
\definecolor{gris65}{gray}{0.35}
\definecolor{gris70}{gray}{0.30}
\definecolor{gris80}{gray}{0.20}
\definecolor{gris85}{gray}{0.15}
\usepackage{graphicx}
\usepackage{float}
\usepackage{caption}
\captionsetup{font=small,labelformat=empty,justification=centering}
\usepackage{booktabs,longtable,array,multirow,calc}
\usepackage{siunitx}
\sisetup{locale=FR}
\usepackage[version=4]{mhchem}
\usepackage{listings}
\usepackage{fancyhdr}
\usepackage{lastpage}
\usepackage[most]{tcolorbox}
\usepackage{enumitem}
\setlist{nosep,leftmargin=*,topsep=2pt}
\usepackage[hidelinks,unicode]{hyperref}
\usepackage{bookmark}
\hypersetup{pdftitle={Agro-Véto TB — Anglais (LVO) 2026},pdfauthor={Rédigé par IA},pdfcreator={LaTeX},pdfsubject={Corrigé guidé}}
\usepackage{xurl}

% --- Pandoc : listes, tableaux, figures ----------------------------------------
\providecommand{\passthrough}[1]{#1}
\providecommand{\tightlist}{\setlength{\itemsep}{0pt}\setlength{\parskip}{0pt}}

\makeatletter
\def\maxwidth{\ifdim\Gin@nat@width>\linewidth\linewidth\else\Gin@nat@width\fi}
\def\maxheight{\ifdim\Gin@nat@height>0.55\textheight 0.55\textheight\else\Gin@nat@height\fi}
\makeatother
\setkeys{Gin}{width=\maxwidth,height=\maxheight,keepaspectratio}
\setcounter{secnumdepth}{-\maxdimen}
\providecommand{\pandocbounded}[1]{#1}
\usepackage{longtable,booktabs,array}

% --- Code (listings) -----------------------------------------------------------
\lstset{
  basicstyle=\ttfamily\small,
  columns=fullflexible,
  keepspaces=true,
  breaklines=true,
  breakatwhitespace=false,
  postbreak=\mbox{\textcolor{gris50}{$\hookrightarrow$}\space},
  frame=single,
  rulecolor=\color{gris35},
  framesep=4pt,
  xleftmargin=5pt,xrightmargin=5pt,
  backgroundcolor=\color{gris3},
  keywordstyle=\bfseries,
  commentstyle=\color{gris65},
  stringstyle=\color{gris80},
  showstringspaces=false,
  tabsize=4,
  extendedchars=true,
  aboveskip=6pt,belowskip=6pt,
}

% --- Macros des corrigés (KaTeX ↔ LaTeX) ---------------------------------------
% Les corrigés sont écrits pour KaTeX : on définit ici ce que KaTeX connaît et
% que LaTeX n'a pas, sans jamais retoucher le Markdown.
\providecommand{\R}{\mathbb{R}}
\providecommand{\N}{\mathbb{N}}
\providecommand{\Z}{\mathbb{Z}}
\providecommand{\Q}{\mathbb{Q}}
\providecommand{\C}{\mathbb{C}}
\providecommand{\dd}{\mathrm{d}}
\providecommand{\llbracket}{\lBrack}
\providecommand{\rrbracket}{\rBrack}
\providecommand{\lBrace}{\lbrace}
\providecommand{\rBrace}{\rbrace}
\providecommand{\Box}{\square}
\providecommand{\Diamond}{\lozenge}
\providecommand{\varnothing}{\emptyset}
\providecommand{\argmax}{\operatorname*{arg\,max}}
\providecommand{\argmin}{\operatorname*{arg\,min}}
\renewcommand{\leq}{\leqslant}
\renewcommand{\geq}{\geqslant}

% Équation hors texte : réduite à la largeur de la ligne si elle déborde.
\newsavebox{\boiteequation}
\newcommand{\equationajustee}[1]{%
  \sbox{\boiteequation}{$\displaystyle #1$}%
  \ifdim\wd\boiteequation>\linewidth
    \[\resizebox{\linewidth}{!}{\usebox{\boiteequation}}\]%
  \else
    \[#1\]%
  \fi}

% --- Encadrés pédagogiques (sobres, lisibles en noir et blanc) -----------------
\tcbset{
  encadre/.style={
    enhanced,breakable,sharp corners,
    colback=white,colframe=gris55,
    boxrule=0.4pt,leftrule=2.2pt,
    left=6pt,right=6pt,top=3pt,bottom=3pt,
    before skip=6pt,after skip=6pt,
    fonttitle=\bfseries\small,coltitle=black,
    attach title to upper,after title={\quad},
    parbox=false,
  }
}
\newtcolorbox{indication}[1]{encadre,colframe=gris45,colback=gris2,title={Indication~#1}}
\newtcolorbox{methode}{encadre,colframe=gris70,title={Méthode}}
\newtcolorbox{piege}{encadre,colframe=gris85,boxrule=0.8pt,leftrule=2.2pt,title={Piège}}
\newtcolorbox{jury}{encadre,colframe=gris60,colback=gris5,title={Jury}}
\newtcolorbox{encadrenote}{encadre,colframe=gris40}

% --- Titres de section ---------------------------------------------------------
\usepackage{titlesec}
\titleformat{\section}{\Large\bfseries}{}{0pt}{}
\titleformat{\subsection}{\large\bfseries}{}{0pt}{}
\titleformat{\subsubsection}{\normalsize\bfseries}{}{0pt}{}
\titlespacing*{\section}{0pt}{14pt plus 4pt minus 2pt}{6pt}
\titlespacing*{\subsection}{0pt}{10pt plus 3pt minus 2pt}{4pt}
\titlespacing*{\subsubsection}{0pt}{8pt plus 2pt minus 2pt}{2pt}

% --- En-tête et pied de page ---------------------------------------------------
\pagestyle{fancy}
\fancyhf{}
\fancyhead[L]{\small\textsc{Corrigé guidé} · Agro-Véto TB}
\fancyhead[R]{\small Anglais 2026}
\fancyfoot[L]{\footnotesize Corrigé rédigé par IA, non officiel}
\fancyfoot[R]{\footnotesize \thepage\,/\,\pageref*{LastPage}}
\renewcommand{\headrulewidth}{0.4pt}
\renewcommand{\footrulewidth}{0pt}
\fancypagestyle{plain}{\fancyhf{}\fancyfoot[R]{\footnotesize \thepage\,/\,\pageref*{LastPage}}\fancyfoot[L]{\footnotesize Corrigé rédigé par IA, non officiel}\renewcommand{\headrulewidth}{0pt}}

\begin{document}
\thispagestyle{plain}
\begin{center}
  {\small\scshape Corrigé guidé}\par\smallskip
  {\LARGE\bfseries Agro-Véto TB — Anglais (LVO) 2026\par}\medskip
  {\large Agro-Véto · TB · Anglais · 2026\par}
\end{center}
\vspace{2pt}\hrule\vspace{8pt}

\emph{Corrigé rédigé par IA, non officiel. Signalez toute erreur.}

\hypertarget{pruxe9sentation-du-sujet}{%
\section{Présentation du sujet}\label{pruxe9sentation-du-sujet}}

\textbf{Format.} L'épreuve dure 2 heures, sans document ni appareil. Elle comprend deux parties rédigées en anglais :

\begin{longtable}[]{@{}
  >{\raggedright\arraybackslash}p{(\columnwidth - 6\tabcolsep) * \real{0.2145}}
  >{\raggedright\arraybackslash}p{(\columnwidth - 6\tabcolsep) * \real{0.4232}}
  >{\raggedright\arraybackslash}p{(\columnwidth - 6\tabcolsep) * \real{0.2106}}
  >{\raggedright\arraybackslash}p{(\columnwidth - 6\tabcolsep) * \real{0.1216}}@{}}
\toprule\noalign{}
\begin{minipage}[b]{\linewidth}\raggedright
Partie
\end{minipage} & \begin{minipage}[b]{\linewidth}\raggedright
Nature
\end{minipage} & \begin{minipage}[b]{\linewidth}\raggedright
Longueur
\end{minipage} & \begin{minipage}[b]{\linewidth}\raggedright
Barème
\end{minipage} \\
\midrule\noalign{}
\endhead
\bottomrule\noalign{}
\endlastfoot
I --- Compréhension de l'écrit & réponse à une question portant sur deux articles, l'un en anglais, l'autre en français, « in your own words » & 220 mots ± 10 \% (198 à 242) & 10 points \\
II --- Expression écrite & argumentation sur l'un des deux sujets, au choix, avec exemples & 220 mots ± 10 \% (198 à 242) & 10 points \\
\end{longtable}

Pour chaque partie, le nombre de mots est indiqué à la fin, et le numéro du sujet choisi en partie II est précisé : le sujet annonce des points de pénalité en cas de non-respect de ces consignes. Ce corrigé traite les deux sujets d'expression, soit trois questions : la compréhension (question 1), le sujet 1 (question 2), le sujet 2 (question 3).

\textbf{Les deux textes.}

\begin{longtable}[]{@{}
  >{\raggedright\arraybackslash}p{(\columnwidth - 4\tabcolsep) * \real{0.3056}}
  >{\raggedright\arraybackslash}p{(\columnwidth - 4\tabcolsep) * \real{0.1454}}
  >{\raggedright\arraybackslash}p{(\columnwidth - 4\tabcolsep) * \real{0.5190}}@{}}
\toprule\noalign{}
\begin{minipage}[b]{\linewidth}\raggedright
Texte
\end{minipage} & \begin{minipage}[b]{\linewidth}\raggedright
Nature
\end{minipage} & \begin{minipage}[b]{\linewidth}\raggedright
Contenu utile
\end{minipage} \\
\midrule\noalign{}
\endhead
\bottomrule\noalign{}
\endlastfoot
Yan Zhuang, « Trump Threatens 200 Percent Tariffs on French Wine », \emph{The New York Times}, 20 janvier 2026 (en anglais) & article d'information politique & menace de droits de douane de 200 \% sur le vin français et le champagne si la France refuse d'entrer au « Board of Peace » ; la menace commerciale comme outil diplomatique ; une menace semblable en mars de l'année précédente, non suivie d'effet, qui avait inquiété les producteurs de champagne, dont les États-Unis sont le premier marché étranger ; réactions politiques françaises et européennes \\
David Aussillou, « Dans le Gard, un rosé ``For sure'' pour parler des taxes américaines », \emph{Ici} (France Bleu), 10 février 2026 (en français) & reportage régional, ton léger & un vigneron des Costières de Nîmes offre au président une bouteille de rosé « For sure » ; près de 200 vignerons vivent de l'appellation et exportent une partie de leur production aux États-Unis ; droits de douane de 10 à 15 \% ; ventes un peu en hausse, mais crainte que la hausse des prix détourne la clientèle américaine ; commandes reçues après le succès viral \\
\end{longtable}

Les deux textes parlent du même conflit, à trois semaines d'écart, mais ne le racontent pas de la même façon. Le premier montre le vin comme \textbf{instrument de pression} dans une crise diplomatique qui le dépasse. Le second montre ses \textbf{effets concrets} sur un vignoble et une \textbf{riposte par la communication}. La question invite à les comparer.

\textbf{Les deux sujets d'expression.} Le sujet 1 généralise la thématique : pourquoi certains produits agricoles deviennent-ils des outils politiques ou symboliques ? Le sujet 2 part de l'idée de l'étiquette « For sure » : la communication et les réseaux sociaux peuvent-ils changer l'impact des tensions politiques ?

\textbf{Ce qu'il ne faut pas laisser passer.} Respecter la fourchette de mots et indiquer un décompte exact, ainsi que le numéro du sujet choisi. En partie I, comparer les deux articles au lieu de les résumer l'un après l'autre, et les nommer précisément. En partie II, définir les termes du sujet et faire servir les exemples à un raisonnement.

\textbf{Gestion du temps.} Une répartition raisonnable : 25 minutes de lecture et de relevé, 30 minutes de rédaction pour la partie I, 45 minutes pour la partie II (choix du sujet, plan au brouillon, rédaction), 15 minutes pour relire la langue et compter les mots. Les deux parties pèsent autant.

\textbf{Ce que dit le rapport 2026.} Le rapport insiste sur la comparaison : « c'est une analyse comparative et non un résumé des textes qui était attendue » (p.~3). Il recense les informations à privilégier et celles à négliger (p.~3), puis pose quatre principes pour l'expression (p.~3-5). Il indique que la première question d'expression « a été très minoritairement choisie » et que la seconde « a été très majoritairement choisie » (p.~3 et 5). Il ne publie ni statistique de notes ni proposition de corrigé complète, mais des « pistes de réflexion » pour le sujet 2 (p.~6).

\hypertarget{partie-i-compruxe9hension-de-luxe9crit}{%
\section{Partie I --- Compréhension de l'écrit}\label{partie-i-compruxe9hension-de-luxe9crit}}

Une seule question, sur deux textes de langue, de ton et de point de vue différents. Le travail consiste à relever les effets des tensions sur le vin, à les classer, puis à les restituer en anglais en comparant les deux approches.

\hypertarget{question-1-how-do-political-tensions-affect-the-wine-sector-according-to-the-two-articles}{%
\subsection{Question 1 --- How do political tensions affect the wine sector according to the two articles?}\label{question-1-how-do-political-tensions-affect-the-wine-sector-according-to-the-two-articles}}

\begin{indication}{1}

Lisez la question mot par mot : \emph{how} (de quelle manière, donc plusieurs effets à distinguer), \emph{affect} (des conséquences, négatives ou non), \emph{the wine sector} (les producteurs et leurs ventes, pas la diplomatie en général), \emph{according to the two articles} (les deux textes, et la façon dont chacun présente le sujet).

\end{indication}

\begin{indication}{2}

Séparez ce qui touche le vin (exportations menacées, incertitude, prix, outil de communication) de ce qui ne le touche pas (Groenland, réactions des responsables politiques, « trade bazooka »). Le second ensemble prend beaucoup de place dans l'article américain, mais ne répond pas à la question.

\end{indication}

\begin{indication}{3}

Pour la comparaison, cherchez le contraste de ton et d'échelle : une menace venue de Washington d'un côté, un vigneron du Gard qui répond par une plaisanterie de l'autre. Annoncez ce contraste dès la première phrase et refermez-le dans la dernière.

\end{indication}

\hypertarget{solution}{%
\subsubsection{Solution}\label{solution}}

Relevé utile, classé par effet :

\begin{longtable}[]{@{}
  >{\raggedright\arraybackslash}p{(\columnwidth - 4\tabcolsep) * \real{0.1868}}
  >{\raggedright\arraybackslash}p{(\columnwidth - 4\tabcolsep) * \real{0.4064}}
  >{\raggedright\arraybackslash}p{(\columnwidth - 4\tabcolsep) * \real{0.3767}}@{}}
\toprule\noalign{}
\begin{minipage}[b]{\linewidth}\raggedright
Effet sur le vin
\end{minipage} & \begin{minipage}[b]{\linewidth}\raggedright
\emph{The New York Times}
\end{minipage} & \begin{minipage}[b]{\linewidth}\raggedright
\emph{Ici} (France Bleu)
\end{minipage} \\
\midrule\noalign{}
\endhead
\bottomrule\noalign{}
\endlastfoot
Le vin, moyen de pression & menace de 200 \% sur le vin et le champagne liée au refus du « Board of Peace » & --- \\
Incertitude & menace de mars de l'année précédente non appliquée, mais qui a inquiété les producteurs de champagne ; les États-Unis, premier marché étranger & les tensions font moins parler d'elles, mais les taxes restent \\
Coût réel & --- & droits de 10 à 15 \% ; près de 200 vignerons concernés ; ventes en légère hausse, mais crainte que la clientèle se détourne \\
Le vin, outil de communication & --- & la cuvée « For sure », le clin d'œil devenu viral, les commandes \\
\end{longtable}

On écarte les éléments que le jury juge « négligeables au regard de cette question » : les réactions politiques, le Groenland, la réaction d'Emmanuel Macron, la riposte européenne.

As both articles show, political tensions affect the wine sector in several ways, although the two texts adopt contrasting angles: The New York Times (20 January 2026) reports a diplomatic threat, while Ici (France Bleu, 10 February 2026) tells a light-hearted story.

In The New York Times, wine appears as a weapon. Donald Trump threatened 200 per cent tariffs on French wine and Champagne if France refused to join his ``Board of Peace''. A similar threat in March 2025 was never carried out, but it worried Champagne producers, whose largest foreign market is the United States. Producers thus live in uncertainty over unrelated disputes.

The France Bleu report shows that the consequences are lasting. The 200 or so winegrowers of the Costières de Nîmes appellation export part of their wine to the United States, where it is now taxed at 10 to 15 per cent. Although their sales there rose slightly last year, professionals fear that higher prices could turn American customers away from European wines.

Yet tensions can also become an opportunity. A winegrower from the Gard named his rosé ``For sure'', after Emmanuel Macron's viral expression at Davos, and offered it to the President, hoping to revive rosé sales in the United States: orders came in immediately.

Both articles thus present wine as a symbolic product caught up in politics: a means of pressure for Washington, but a means of communication for French producers.

\textbf{(236 words)}

Choix justifiés :

\begin{itemize}
\tightlist
\item
  La phrase d'ouverture reprend presque mot pour mot la formulation que suggère le jury, puis annonce le contraste entre les deux articles.
\item
  Chaque article est nommé par son titre de presse (\emph{The New York Times}, \emph{France Bleu}), et non par « the first article » ou « the French article ».
\item
  « March last year » est rendu par \emph{March 2025}, puisque l'article date de janvier 2026.
\item
  La dernière phrase compare les deux rôles du vin, moyen de pression et moyen de communication : c'est la mise en relation que le jury demande.
\end{itemize}

\begin{methode}

Pour une question comparative, organisez la réponse par idées, pas par article : chaque paragraphe dit un effet et indique dans quel texte on le trouve. Les phrases d'ouverture et de conclusion portent la comparaison (\emph{although the two texts adopt contrasting angles}, \emph{both articles thus present\ldots{}}).

\end{methode}

\begin{jury}

Rapport 2026 (p.~2) : les différences et les similitudes entre les textes « devaient absolument être soulignées » (« menace diplomatique d'un côté, humour de l'autre »), et « L'aspect comparatif de la question devait donc transparaître à la lecture de la production du/de la candidat(e). » Le jury suggère d'ouvrir ainsi : « Political tensions affect the wine sector in several/different ways, as shown by both articles\ldots{} » (p.~3).

\end{jury}

\begin{jury}

Rapport 2026 (p.~3) : les bonnes copies ont retenu ce qui touchait au vin (« exportations menacées, incertitude économique, le vin comme outil de communication, de politique, symbolique\ldots{} ») au détriment d'éléments négligeables (« les réactions du monde politique, la question du Groenland, la réaction de M. Macron, la riposte de l'Europe\ldots{} »).

\end{jury}

\begin{piege}

Les sources doivent être identifiables. Pour le jury, renvoyer aux textes par « the first article/the second article » ou « the French article/the English article » « n'est pas suffisant », pas plus que par le seul nom des auteurs (p.~3). Nommez le journal, une fois au moins avec la date.

\end{piege}

\begin{piege}

Le texte français ne se recopie pas : « droits de douane » se dit \emph{tariffs} ou \emph{customs duties}, « vigneron » \emph{winegrower} ou \emph{winemaker}, « clin d'œil » \emph{nod} ou \emph{wink}, « millésime » \emph{vintage}. Le jury rappelle que l'on peut, faute de mieux, reprendre certains termes du texte anglais, « ce qui n'est pas possible à partir d'un texte en français » (p.~2).

\end{piege}

\hypertarget{partie-ii-expression-uxe9crite}{%
\section{Partie II --- Expression écrite}\label{partie-ii-expression-uxe9crite}}

Le candidat choisit un seul des deux sujets. Le premier demande de l'analyse et des exemples historiques ; le second, choisi par la majorité des candidats, demande de démontrer un lien de cause à effet entre communication et tensions. Les deux devoirs ci-dessous suivent le même schéma : contexte et question, trois paragraphes argumentés, conclusion qui répond.

\hypertarget{question-2-sujet-1-why-do-certain-agricultural-products-such-as-wine-wheat-or-dairy-products-become-political-or-symbolic-tools-in-international-relations-give-examples.}{%
\subsection{Question 2 --- Sujet 1 : Why do certain agricultural products (such as wine, wheat, or dairy products) become political or symbolic tools in international relations? Give examples.}\label{question-2-sujet-1-why-do-certain-agricultural-products-such-as-wine-wheat-or-dairy-products-become-political-or-symbolic-tools-in-international-relations-give-examples.}}

\begin{indication}{1}

La question commence par \emph{why} : on attend des \textbf{causes}, pas une liste d'exemples. Chaque paragraphe doit donc répondre « parce que\ldots{} », et l'exemple vient ensuite le prouver.

\end{indication}

\begin{indication}{2}

Pesez chaque terme : \emph{certain} (tous les produits n'ont pas la même valeur), \emph{political} \textbf{et} \emph{symbolic} (deux dimensions à traiter), \emph{agricultural products} (ni le pétrole ni le gaz), \emph{international relations} (diplomatie, commerce, conflits).

\end{indication}

\begin{indication}{3}

Trois raisons possibles : certains produits sont vitaux (le blé), d'autres incarnent une identité nationale (le champagne, les fromages AOP), et leurs producteurs pèsent sur les gouvernements. Pour chacune, prévoyez un exemple daté que vous maîtrisez.

\end{indication}

\hypertarget{solution-1}{%
\subsubsection{Solution}\label{solution-1}}

On suit ce plan : une introduction qui part de la menace américaine sur le vin français, puis trois causes, chacune prouvée par un exemple (le blé ukrainien bloqué en mer Noire, le roquefort et le vin taxés par les États-Unis, le vin australien taxé par la Chine), et une conclusion qui répond et nuance.

\textbf{Subject 1}

In January 2026, Donald Trump threatened to tax French wine at 200 per cent because France had refused to join his ``Board of Peace''. Why do certain agricultural products become political or symbolic tools in international relations?

Not all products carry the same weight. First, some are vital, which makes them powerful levers. Wheat feeds whole populations, so controlling its supply gives a country leverage. After invading Ukraine in 2022, Russia blocked Ukrainian grain exports through the Black Sea, and food prices rose in many importing countries. Food thus becomes a weapon.

Secondly, other products are symbols of national identity. Champagne, Roquefort or Scotch whisky embody a country's know-how and are protected by designations of origin. Hitting them hurts a nation's pride as much as its economy. This is why the United States imposed a 100 per cent duty on Roquefort in 1999 during the dispute over hormone-treated beef, and taxed French wine in 2019 during the Airbus-Boeing conflict.

Thirdly, agricultural products affect voters. Farmers are well organised in many countries, so targeting their exports puts indirect pressure on governments. When China taxed Australian wine heavily in 2020, after Canberra had called for an inquiry into the origins of Covid-19, Australian winemakers lost their biggest export market.

Agricultural products therefore become tools because they combine economic importance, emotional value and political weight. Yet such pressure seldom changes a government's policy and mainly hurts producers.

\textbf{(234 words)}

Choix justifiés :

\begin{itemize}
\tightlist
\item
  La première phrase du développement (\emph{Not all products carry the same weight}) répond au mot \emph{certain} du sujet.
\item
  Le deuxième paragraphe définit la dimension \textbf{symbolique} : le produit incarne un savoir-faire et une fierté nationale, protégés par une appellation.
\item
  Le troisième traite de la \textbf{pression indirecte} sur les gouvernements par l'intermédiaire des agriculteurs, que le jury dit « largement ignoré » par les candidats.
\item
  Les exemples sont datés et sûrs : droits de 100 \% sur le roquefort en 1999 (conflit du bœuf aux hormones), taxation du vin français en 2019 (conflit Airbus-Boeing), droits chinois sur le vin australien en 2020.
\end{itemize}

\begin{methode}

Une question en \emph{why} se traite par causes : \emph{First, some are vital\ldots{}}, \emph{Secondly, other products are symbols\ldots{}}, \emph{Thirdly, agricultural products affect voters\ldots{}}. Chaque phrase-thèse doit pouvoir se lire comme une réponse à la question, même sans son exemple.

\end{methode}

\begin{jury}

Rapport 2026 (p.~4) : si \emph{political tool} a été bien analysé, « l'expression symbolical tool n'a pas toujours été clairement définie et développée ». Le jury ajoute que les produits agricoles « pouvaient être présentés comme un moyen de pression indirecte sur le débat public et sur les décisions politiques. Cet aspect a été largement ignoré par les candidat(e)s. »

\end{jury}

\begin{piege}

Le jury relève des exemples hors sujet ou erronés : le gaz naturel et le pétrole, qui « ne sont pas considérés comme des produits agricoles » (p.~4), et d'autres exemples « farfelus ou erronés » (p.~5). Un exemple dont on n'est pas sûr est plus dangereux qu'utile.

\end{piege}

\begin{jury}

Rapport 2026 (p.~4) : le jury met en garde contre les reformulations de la question en fin d'introduction, souvent fautives (\emph{*How the international tradding can become a political tool?}), et conclut que la question « ne nécessitait pas en réalité de reformulation particulière ». Le devoir proposé la pose donc telle quelle.

\end{jury}

\hypertarget{question-3-sujet-2-do-you-think-social-media-and-communication-strategies-can-change-the-impact-of-political-tensions-give-examples.}{%
\subsection{Question 3 --- Sujet 2 : Do you think social media and communication strategies can change the impact of political tensions? Give examples.}\label{question-3-sujet-2-do-you-think-social-media-and-communication-strategies-can-change-the-impact-of-political-tensions-give-examples.}}

\begin{indication}{1}

Définissez les termes avant d'écrire : \emph{political tensions} (conflits diplomatiques ou commerciaux), \emph{impact} (les conséquences concrètes : ventes, opinion, décisions), \emph{social media} et \emph{communication strategies} (deux choses différentes : un canal de diffusion, une stratégie volontaire).

\end{indication}

\begin{indication}{2}

\emph{Change} peut aller dans les deux sens : les réseaux sociaux peuvent \textbf{aggraver} l'impact des tensions comme l'\textbf{atténuer}. Un plan qui examine les deux, puis fixe les limites, répond à la question en entier.

\end{indication}

\begin{indication}{3}

Gardez une place pour la limite : une étiquette virale ne supprime pas un droit de douane. C'est elle qui permet de mesurer la portée du mot \emph{impact} et de conclure avec nuance.

\end{indication}

\hypertarget{solution-2}{%
\subsubsection{Solution}\label{solution-2}}

On suit ce plan : une introduction qui part de la cuvée « For sure » et pose la question, une première partie sur l'amplification des tensions, une deuxième sur la communication qui atténue ou retourne leurs effets (la cuvée, l'Ukraine, le Canada), une troisième sur les limites, puis une conclusion qui donne un avis personnel.

\textbf{Subject 2}

When Donald Trump threatened French wine with 200 per cent tariffs, a winegrower from the Gard replied with a rosé called ``For sure'', inspired by Emmanuel Macron's viral words at Davos. Can social media and communication strategies change the impact of political tensions?

To begin with, social media amplify tensions. Political leaders now announce sanctions or tariffs directly online, as Donald Trump often does on Truth Social. Threats therefore spread instantly and worry markets and producers before any decision has been taken. A diplomatic quarrel can quickly turn into a boycott.

Nevertheless, communication can also soften the effects of a crisis, or even turn it into an opportunity. The ``For sure'' label made a tariff dispute amusing and immediately brought in orders. Similarly, in 2022, Volodymyr Zelensky's daily videos helped Ukraine win international sympathy and military aid. In Canada, the ``Buy Canadian'' movement that spread online in 2025 in response to American tariffs encouraged consumers to support local producers.

However, the influence of these strategies remains limited. A viral bottle does not remove customs duties: Costières wines still pay 10 to 15 per cent to enter the United States. Communication changes the way tensions are perceived, not the political decisions themselves.

In conclusion, I believe that social media and communication strategies can change the impact of political tensions, mainly by shaping perceptions and creating new markets, but they cannot replace diplomacy.

\textbf{(231 words)}

Choix justifiés :

\begin{itemize}
\tightlist
\item
  Chaque exemple est relié à l'idée qu'il démontre : la publication de menaces en ligne prouve l'amplification ; la cuvée et le mouvement « Buy Canadian » prouvent que la communication peut retourner une crise.
\item
  La limite reprend un fait du texte de France Bleu (les droits de 10 à 15 \% restent dus) : l'article peut nourrir l'argumentation.
\item
  La conclusion commence par \emph{I believe}, puisque le sujet demande un avis (\emph{Do you think}).
\end{itemize}

\begin{methode}

Pour démontrer un lien de cause à effet, enchaînez trois étapes : le mécanisme (\emph{threats spread instantly}), sa conséquence (\emph{worry markets and producers}), l'exemple qui le prouve. Un exemple posé sans mécanisme n'est qu'une illustration.

\end{methode}

\begin{jury}

Rapport 2026 (p.~5) : cette question « a été très majoritairement choisie », mais « Bien souvent le lien entre les tensions politiques et les réseaux sociaux, les stratégies de communication n'a pas été démontré. La portée du terme « impact » n'a pas non plus été mesurée. » Le jury ajoute qu'une liste d'avantages et d'inconvénients des réseaux sociaux « ne prouve rien en soi » (p.~5-6).

\end{jury}

\begin{jury}

Rapport 2026 (p.~6) : parmi les pistes proposées, la limite : « social media and communication strategies can influence perception (how tensions are perceived) but their impact remains limited ». Le candidat a « parfaitement le droit d'exprimer une opinion personnelle. Celle-ci en revanche, doit être argumentée. »

\end{jury}

\begin{piege}

\emph{Media} est déjà un pluriel : on écrit \emph{social media} sans \emph{s}, et le jury relève \emph{*social medias are} (p.~6). De même, \emph{information} et \emph{news} sont indénombrables (\emph{*informations}), et \emph{the United States} s'accorde au singulier. La question directe se construit avec l'auxiliaire : \emph{Can social media change\ldots?}, et non \emph{*Social media can change \ldots{} ?} (p.~7).

\end{piege}

\hypertarget{bilan}{%
\section{Bilan}\label{bilan}}

\textbf{À retenir pour la partie I.}

\begin{itemize}
\tightlist
\item
  Lire la question mot par mot : ici, \emph{according to the two articles} impose une comparaison.
\item
  Relever, trier (ce qui touche le vin, ce qui ne le touche pas), classer par effet.
\item
  Organiser par idées, pas par texte ; ouvrir et fermer sur la comparaison.
\item
  Nommer chaque source par son titre de presse ; ne pas se contenter de « the French article ».
\end{itemize}

\textbf{À retenir pour la partie II.}

\begin{itemize}
\tightlist
\item
  Contextualiser de façon pertinente, sans généralité vague ; poser la question sans la déformer.
\item
  Définir les termes clés (\emph{symbolic}, \emph{impact}), qui orientent tout le plan.
\item
  Faire précéder chaque exemple de l'idée qu'il démontre ; conclure par une réponse nette et nuancée.
\end{itemize}

\textbf{Le décompte des mots.} Il doit être indiqué, tout comme le numéro du sujet choisi, et le jury précise « qu'un comptage manifestement insincère sera sanctionné » (p.~7). Viser 225 à 235 mots laisse une marge si les conventions de décompte diffèrent (traits d'union, nombres).

\textbf{Langue à revoir.} Le jury recommande de travailler la détermination, les pronoms, la place de l'adjectif, les temps, les verbes irréguliers, la voix passive (« source de nombreuses erreurs cette année encore », p.~7), l'expression de la durée (\emph{since}, \emph{for}), les modaux et la construction des questions. Lexique de base à maîtriser : \emph{threat} / \emph{to threaten}, \emph{tariff}, \emph{policy} / \emph{politics} / \emph{politician}, \emph{produce} / \emph{producer}.

\textbf{Situations qui se retrouvent dans d'autres sujets.} Comparer un article d'information et un reportage de ton différent sur un même événement. Expliquer pourquoi un bien économique devient un enjeu politique (alimentation, énergie, technologie). Évaluer le pouvoir réel de la communication face à une décision politique.

\end{document}
